\documentclass[conference,compsoc]{IEEEtran}

\usepackage{amsmath,amsfonts,amssymb}
\usepackage{algorithmic}
\usepackage{algorithm}
\usepackage{array}
\usepackage[caption=false,font=footnotesize,labelfont=sf,textfont=sf]{subfig}
\usepackage{textcomp}
\usepackage{stfloats}
\usepackage{url}
\usepackage{verbatim}
\usepackage{graphicx}
\usepackage{listings}
\usepackage{xcolor}
\usepackage[catalan]{babel}
\usepackage[T1]{fontenc}
\usepackage[utf8]{inputenc}
\usepackage{booktabs}
\usepackage{float}
\usepackage{cuted}
\usepackage{microtype}
\usepackage{pgfplots}
\pgfplotsset{compat=1.18}
\graphicspath{{./}}

\renewcommand{\abstractname}{Resum}
\renewcommand{\IEEEkeywordsname}{Termes d'índex}

\makeatletter
\def\abstract{\normalfont\@IEEEtweakunitybaselinestretch{1.15}%
    \if@twocolumn
      \@IEEEabskeysecsize\noindent\textbf{\textit{\abstractname}}---\normalfont\@IEEEabskeysecsize\relax
    \else
      \bgroup\par\addvspace{0.5\baselineskip}\centering\vspace{-1.78ex}\@IEEEabskeysecsize\textbf{\abstractname}\par\addvspace{0.5\baselineskip}\egroup\quotation\normalfont\@IEEEabskeysecsize%
    \fi\@IEEEgobbleleadPARNLSP}
\def\IEEEkeywords{\normalfont\@IEEEtweakunitybaselinestretch{1.15}%
    \if@twocolumn
      \@IEEEabskeysecsize\vskip 0.5\baselineskip plus 0.25\baselineskip minus 0.25\baselineskip\noindent
      \textbf{\textit{\IEEEkeywordsname}}---\normalfont\@IEEEabskeysecsize\relax
    \else
      \bgroup\par\addvspace{0.5\baselineskip}\centering\@IEEEabskeysecsize\textbf{\IEEEkeywordsname}\par\addvspace{0.5\baselineskip}\egroup\quotation\normalfont\@IEEEabskeysecsize%
    \fi\@IEEEgobbleleadPARNLSP}
\makeatother

\lstset{
    language=Java,
    basicstyle=\ttfamily\footnotesize,
    keywordstyle=\color{blue}\bfseries,
    commentstyle=\color{gray}\itshape,
    stringstyle=\color{red!70!black},
    numbers=left,
    numberstyle=\tiny\color{gray},
    numbersep=5pt,
    frame=single,
    breaklines=true,
    showstringspaces=false,
    tabsize=4,
    captionpos=b
}

\hyphenation{Mon-te-Car-lo per-cen-til per-cen-tils si-mu-la-cio con-tro-la-dor Swing-Worker re-pro-du-i-ble}

\begin{document}

\title{Simulació Probabilística del Joc de l'Oca\\Pràctica 7}

\author{\IEEEauthorblockN{Serafí Nebot Ginard}
\IEEEauthorblockA{Universitat de les Illes Balears (UIB)\\
Grau en Enginyeria Informàtica\\
Algorismes Avançats, Curs 2025--2026}}

\maketitle

\begin{strip}
\begin{abstract}
Aquesta memòria descriu una aplicació Java per estimar mitjançant simulació Monte
Carlo el nombre de torns necessaris per acabar una partida del Joc de l'Oca. La
implementació reprodueix les regles exactes de l'enunciat per a un jugador i
afegeix extensions opcionals per a 2, 3 i 4 jugadors amb captura de fitxes,
comparació de variants, freqüències de visita i execució concurrent. El projecte s'ha organitzat amb
arquitectura Model-Vista-Controlador, una entrada de consola obligatòria, una
interfície Swing i una eina CLI específica per generar experiments reproduïbles per
a la memòria. La part central del treball és l'anàlisi empírica: s'estudia la
convergència de la mitjana, l'estabilitat dels percentils, la cua de la
distribució, l'efecte de variants de regles, les probabilitats de victòria segons
el nombre de jugadors, les freqüències de visita de caselles i l'impacte de la
simulació concurrent. Amb $250\,000$ partides per repetició en les comparacions
principals, la mitjana d'una
partida individual queda al voltant de $20.75$ torns, amb mediana 17 i percentil
$0.9$ al voltant de 38--39 torns.
\end{abstract}

\begin{IEEEkeywords}
Monte Carlo, algorismes probabilístics, simulació, percentils empírics, Joc de
l'Oca, multijugador, MVC, Java Swing, experiments reproduïbles.
\end{IEEEkeywords}
\end{strip}

\IEEEpeerreviewmaketitle

%% ================================================================
\section{Introducció}

La pràctica demana estimar el nombre de torns necessaris per finalitzar una
partida del Joc de l'Oca amb un únic jugador. El tauler té 64 caselles numerades
de 0 a 63, la partida comença a la casella 0 i acaba immediatament quan el jugador
arriba a la casella 63. En cada torn ordinari es llança un dau equilibrat de sis
cares. Després del moviment normal i del possible rebot respecte de la casella
final, es comprova si la casella assolida és especial. Com a ampliació opcional,
la mateixa implementació permet simular partides de 1 a 4 jugadors i estimar la
probabilitat de victòria de cada posició d'ordre de torn. En aquest mode, si una
fitxa acaba el moviment en una casella ocupada per una altra fitxa, la fitxa que hi
era abans torna a la casella inicial.

El problema és probabilístic perquè el resultat de cada tirada és aleatori. Es
podria construir un model exacte com una cadena de Markov finita, amb estats que
incloguessin la posició i les penalitzacions pendents. En aquesta pràctica, però,
l'objectiu és aplicar simulació Monte Carlo: generar moltes partides independents,
mesurar quants torns dura cadascuna i resumir la mostra amb mínim, percentils,
mediana, màxim i mitjana.

La implementació no es limita al format de consola. S'ha construït una aplicació
amb tres entrades:
\begin{itemize}
    \item \texttt{Main}, que compleix estrictament el format obligatori de
    l'enunciat;
    \item \texttt{GUI}, que obre una interfície Swing per executar simulacions,
    comparar variants, triar el nombre de jugadors i visualitzar la distribució;
    \item \texttt{CLI}, que executa experiments deterministes amb llavors fixes,
    repeticions, variants, nombre de jugadors i sortida Markdown o CSV.
\end{itemize}

A aquesta memòria primer es
formalitza el problema, després es descriu el disseny de l'aplicació i finalment
es dedica una part extensa als experiments. En particular, la Secció
\ref{sec:experiments} és la part més important: totes les taules provenen de la
CLI reproduïble i no de proves manuals de la GUI.

%% ================================================================
\section{Descripció del Problema}

\subsection{Estat de la Partida}

En el cas obligatori d'un jugador, l'estat necessari per continuar una partida és
molt petit:
\begin{equation}
S=(p,a)
\end{equation}
on $p\in\{0,\ldots,63\}$ és la posició actual i $a\in\{0,1,2,3\}$ és el nombre de
torns anul·lats pendents. El valor màxim de $a$ és 3 perquè la penalització més
llarga de l'enunciat és la presó. La casella 63 és absorbent: un cop assolida, la
partida s'acaba i no es processen més efectes.

La variable aleatòria que s'estima és $X$, definida com el nombre de torns
comptats fins que la partida arriba a la casella 63.
Un torn anul·lat per penalització compta dins $X$, però no llança el dau ni mou el
jugador. En canvi, les tirades addicionals d'una oca pertanyen al mateix torn i no
incrementen el comptador.

Per a l'ampliació multijugador, cada jugador té el seu propi parell $(p,a)$ i el
torn avança cíclicament. La partida acaba en el moment que el jugador actual arriba
a la casella 63. Per tant, si hi ha $J$ jugadors, l'estat global és
\begin{equation}
S=((p_1,a_1),\ldots,(p_J,a_J),j)
\end{equation}
on $j$ indica quin jugador ha de moure. La variable d'interès continua essent el
nombre total de torns comptats fins al final de la partida, però també es registra
el guanyador $W\in\{1,\ldots,J\}$. Després de cada torn multijugador, si la posició
final del jugador actual coincideix amb la d'una altra fitxa, l'ocupant anterior es
reinicia a la casella 0 i perd qualsevol penalització pendent. Aquesta captura és
l'única interacció directa entre jugadors.

\subsection{Moviment i Rebot}

Si el jugador és a la casella $p$ i obté un valor de dau $d$, primer es calcula
$r=p+d$. Si $r\leq63$, la casella provisional és $r$. Si $r>63$, s'aplica el rebot
definit a l'enunciat:
\begin{equation}
\label{eq:rebot}
q = 63 - (r-63).
\end{equation}
Per exemple, des de la casella 60 amb un 5, $r=65$ i la posició resultant és
$61$. Aquest rebot fa que acostar-se a la meta no sigui equivalent a acabar: una
tirada massa gran pot allunyar el jugador.

\subsection{Caselles Especials}

Després del moviment ordinari, i només després, s'aplica com a màxim un efecte
especial. La Taula~\ref{tab:regles} resumeix les regles implementades. La regla
d'un únic efecte és essencial: si una tirada porta a la casella 6, el jugador passa
a la 12, però no s'aplica immediatament l'efecte de la 12. La implementació
conserva internament tant la casella de caiguda com la casella final del moviment
per poder comptar visites sense encadenar efectes.

\begin{table}[H]
\centering
\caption{Caselles especials implementades segons l'enunciat.}
\label{tab:regles}
\footnotesize
\begin{tabular}{@{}p{0.26\columnwidth}p{0.68\columnwidth}@{}}
\toprule
Tipus & Efecte \\
\midrule
Oques & $5\to9\to14\to18\to23\to27\to32\to36\to41\to45\to50\to54\to59\to63$; dona una nova tirada sense sumar torn si no acaba. \\
Ponts & $6\leftrightarrow12$. \\
Dados & $26\leftrightarrow53$. \\
Posada & Casella 19; un torn anul·lat. \\
Pou & Casella 31; dos torns anul·lats. \\
Laberint & $42\to30$. \\
Presó & Casella 52; tres torns anul·lats. \\
Mort & $58\to0$. \\
\bottomrule
\end{tabular}
\end{table}

\subsection{Variants Analitzades}

A més de les regles oficials, el programa permet variants controlades. No formen
part de la sortida obligatòria, però són útils per entendre quin efecte té cada
família de caselles sobre la distribució de torns. Les variants són:
\begin{itemize}
    \item \texttt{STANDARD}: totes les regles oficials.
    \item \texttt{NO\_PENALTIES}: manté oques, ponts, dados, laberint i mort, però
    elimina posada, pou i presó.
    \item \texttt{NO\_GOOSE\_EXTRA\_ROLL}: manté els salts d'oca, però la tirada
    addicional passa a comptar com un altre torn.
    \item \texttt{NO\_SPECIAL\_CELLS}: desactiva totes les caselles especials i
    deixa només el moviment amb rebot.
\end{itemize}

%% ================================================================
\section{Fonaments Teòrics}

\subsection{Estimació Monte Carlo}

Una simulació de $N$ partides genera observacions independents
$x_1,x_2,\ldots,x_N$ de la variable aleatòria $X$. La mitjana mostral és
\begin{equation}
\label{eq:mitjana}
\bar{x}=\frac{1}{N}\sum_{i=1}^{N}x_i.
\end{equation}
Per la llei dels grans nombres, $\bar{x}$ convergeix cap a $E[X]$ quan $N$ creix.
Si les partides són independents i $X$ té variància finita $\sigma^2$, llavors
\begin{equation}
\label{eq:error-mitjana}
\operatorname{Var}(\bar{x}) = \frac{\sigma^2}{N},
\end{equation}
de manera que la desviació típica de l'estimador és $\sigma/\sqrt{N}$.
Aquest és el motiu pel qual s'espera que l'error típic
de la mitjana decreixi aproximadament com $O(1/\sqrt{N})$. Això no significa que
qualsevol estadístic sigui igual d'estable. La mitjana i els percentils centrals
convergeixen ràpidament, però el màxim observat tendeix a créixer quan augmenta
$N$, perquè depèn de la cua de la distribució.

\subsection{Percentils Empírics}

L'enunciat demana els percentils empírics 0.1, 0.2, 0.3, 0.4, la mediana, 0.6,
0.7, 0.8 i 0.9. La implementació utilitza el criteri de rang més proper. Si
$y_1\leq y_2\leq\cdots\leq y_N$ és la mostra ordenada, llavors
\begin{equation}
\label{eq:percentil}
Q_p = y_{\lceil pN \rceil}.
\end{equation}
Aquest criteri té dos avantatges pràctics: sempre retorna un nombre enter de torns
realment observat i és fàcil de reproduir a partir de la mostra ordenada. La
mitjana és l'únic valor de sortida que pot no ser enter i s'imprimeix amb sis
decimals.

\subsection{Cost Asimptòtic}

Si $T$ és el nombre mitjà de torns d'una partida, simular $N$ partides té cost
esperat $O(NT)$. El càlcul de percentils ordena els $N$ valors observats i costa
$O(N\log N)$. Per tant, el cost total de la simulació completa amb estadístics és
\begin{equation}
\label{eq:cost}
O(NT + N\log N).
\end{equation}
La memòria és $O(N+64+J)$: cal guardar els $N$ torns per ordenar-los, un vector de
64 posicions per comptar visites a caselles i un vector de $J$ comptadors de
victòries quan s'executen partides multijugador.

Amb $h$ fils, la part de simulació es reparteix en blocs independents. Idealment,
el terme $NT$ es divideix per $h$, però l'ordenació final i la sobrecàrrega de
crear fils, combinar vectors i gestionar memòria no desapareixen. Per això no
s'espera un guany lineal perfecte.

%% ================================================================
\section{Disseny i Implementació}

\subsection{Estructura del Projecte}

La separació MVC queda reflectida en tres paquets:
\texttt{model}, \texttt{controller} i \texttt{view}. Les entrades principals són
\texttt{Main}, \texttt{GUI} i \texttt{CLI}. La Taula~\ref{tab:classes} resumeix les
classes més importants.

\begin{table}[H]
\centering
\caption{Responsabilitats principals de la implementació.}
\label{tab:classes}
\scriptsize
\begin{tabular}{@{}p{0.48\columnwidth}p{0.46\columnwidth}@{}}
\toprule
Classe & Responsabilitat \\
\midrule
\texttt{Main} & Entrada obligatòria de consola amb el format exacte de l'enunciat. \\
\texttt{GUI} & Entrada gràfica; només crea la finestra i connecta vista i controlador. \\
\texttt{CLI} & Experiments reproduïbles per a la memòria. \\
\texttt{Game} & Estat i regles d'una partida de 1--4 jugadors, efectes especials i tirades del dau. \\
\texttt{GameSnapshot} & Posició i penalitzacions pendents entre torns. \\
\texttt{MoveResult} & Resultat intern d'una tirada: casella de caiguda, casella final i penalització. \\
\texttt{MatchResult} & Durada i guanyador d'una partida completa. \\
\texttt{Simulation} & Execució de moltes partides, mode seqüencial i mode concurrent. \\
\texttt{SimulationData} & Dades brutes: torns, visites i victòries. \\
\texttt{SimulationStats} & Ordenació de mostres, percentils, mínim, màxim i mitjana. \\
\texttt{SimulationResult} & Resultat preparat per a la GUI. \\
\texttt{RuleVariant} & Configuració de variants de regles. \\
\texttt{SimulationController} & Connexió entre GUI i model; llança càlculs amb \texttt{SwingWorker}. \\
\texttt{SimulationView} & Contracte mínim que ha de complir la vista gràfica. \\
\texttt{SimulationViewListener} & Esdeveniments de la vista que el controlador ha d'atendre. \\
\texttt{SimulationFrame} & Formulari, taules, sortida requerida i pestanyes de resultats. \\
\texttt{HistogramPanel} & Representació gràfica discreta de la distribució de torns. \\
\bottomrule
\end{tabular}
\end{table}

\subsection{Entrades del Programa}

El projecte té tres punts d'entrada perquè cada ús té una responsabilitat diferent.
\texttt{Main} és el programa obligatori de l'enunciat: llegeix un únic enter $N$,
executa $N$ partides amb regles oficials i imprimeix exactament mínim, percentils,
mediana, màxim i mitjana. No accepta llavors ni opcions addicionals per evitar que
la sortida requerida quedi contaminada amb text de depuració o de benchmark.

\texttt{GUI} és l'entrada gràfica. Inicialitza Swing, crea la finestra principal
i la connecta amb el controlador. Això manté la classe d'arrencada separada tant
de la lògica de simulació com de la construcció visual detallada.

\texttt{CLI} és l'entrada experimental. A diferència de \texttt{Main}, sí accepta
paràmetres com $N$, llavor, variants, repeticions, nombre de jugadors, fils, format
CSV o Markdown i nombre de caselles visitades a mostrar. Aquesta separació permet
tenir alhora una sortida estricta per a la correcció automàtica i una eina flexible
per reconstruir les taules de la memòria.

\subsection{Model del Joc}

La classe \texttt{Game} és el centre del model. Conté les constants del tauler, les
caselles especials, la tirada del dau, el rebot, la regla d'un únic efecte especial,
les tirades extra d'oca, les penalitzacions, la mort, el laberint i les captures
multijugador. En el cas d'un jugador, la lògica principal és la següent:

\begin{lstlisting}[caption={Esquema de simulació d'una partida.},label={lst:partida}]
int turns = 0;
GameSnapshot state = new GameSnapshot(0, 0);
while (!state.isFinished()) {
    state = advanceTurn(state, visits);
    turns++;
}
\end{lstlisting}

El mètode \texttt{advanceTurn} distingeix dos casos. Si hi ha penalitzacions
pendents, decrementa el comptador de penalització i retorna el mateix lloc. Si no
n'hi ha, llança el dau i aplica \texttt{applyRoll}. Quan una oca dona una tirada
extra, el mètode repeteix el procés dins el mateix torn comptat. Per això el mínim
observat pot arribar a ser 1: una cadena favorable d'oques pot portar directament
fins a la meta sense incrementar el comptador de torns.

\texttt{GameSnapshot} és un registre immutable amb la posició i les penalitzacions
pendents d'una fitxa entre dos torns. S'utilitza també en multijugador: una partida
amb $J$ jugadors manté un vector de $J$ snapshots. La immutabilitat evita efectes
laterals difícils de seguir, perquè avançar un torn sempre produeix un nou estat.

\texttt{MoveResult} és el resultat intern d'una tirada de dau. Guarda la casella on
s'ha caigut, la casella final després de l'efecte especial, les penalitzacions
generades, si hi ha una tirada extra i si la partida ha acabat. Aquesta classe és
necessària perquè una tirada pot tenir dues caselles rellevants: per exemple, caure
a 6 i acabar a 12 pel pont.

\texttt{MatchResult} representa una partida completa. Conté el nombre total de torns
i el guanyador. En el cas obligatori d'un jugador, el guanyador sempre és el jugador
0; en l'ampliació multijugador, aquest camp alimenta les probabilitats de victòria.

\texttt{RuleVariant} encapsula quins grups de regles estan actius. Les variants no
afecten \texttt{Main}, però permeten que la CLI i la GUI comparin regles oficials,
absència de penalitzacions, oques sense tirada extra i absència de caselles
especials sense duplicar el codi de moviment.

En mode multijugador, \texttt{Game} manté un vector de \texttt{GameSnapshot}, un per
jugador, i un índex de torn. Cada torn només modifica primer l'estat del jugador
actual. Si aquest jugador arriba a la casella 63, es retorna un \texttt{MatchResult}
amb el nombre total de torns i l'índex del guanyador. Si no acaba i la seva casella
final ja estava ocupada, l'ocupant anterior es reinicia a la casella 0 i perd les
penalitzacions pendents. Aquesta regla de captura s'aplica després d'haver resolt el
moviment, el rebot i com a màxim un efecte especial.

\subsection{Simulació i Estadística}

La classe \texttt{Simulation} executa moltes partides independents. En mode
seqüencial, totes les partides reutilitzen un \texttt{Game} amb un generador
\texttt{Random}. En mode concurrent, el nombre total de partides es divideix en
blocs. Cada bloc rep una llavor derivada de la llavor base i de l'índex del
treballador, de manera que no hi ha cap \texttt{Random} compartit entre fils.

El resultat brut és un \texttt{SimulationData}: un vector amb els torns de cada
partida, un vector de 64 comptadors de visita i un vector de comptadors de
victòries per jugador. Aquesta classe només transporta observacions i clona els
vectors d'entrada i sortida per evitar modificacions accidentals.

\texttt{SimulationStats} rep els torns bruts i en calcula mínim, màxim, mitjana,
mediana i percentils. L'ordenació de la mostra queda concentrada aquí, de manera que
\texttt{Main}, \texttt{CLI} i la GUI comparteixen exactament el mateix criteri
estadístic.

\texttt{SimulationResult} és l'objecte de resultat de la GUI. A més dels estadístics,
manté els torns originals per dibuixar l'histograma, les visites per casella, els
comptadors de victòria, la llavor usada, el nombre de fils i, si escau, la
comparació de variants.

\subsection{Controlador}

\texttt{SimulationController} connecta la vista amb el model. Rep els paràmetres
introduïts per l'usuari, valida el nombre de partides, decideix quines variants cal
executar i construeix el \texttt{SimulationResult}. Les simulacions s'executen dins
un \texttt{SwingWorker} per no bloquejar el fil d'esdeveniments de Swing. Així,
la finestra pot actualitzar l'estat d'execució i mostrar errors sense conèixer els
detalls de la simulació.

\subsection{Vista Gràfica}

\texttt{SimulationView} és la interfície que el controlador utilitza per comunicar-se
amb qualsevol vista. Defineix operacions com canviar l'estat d'execució, mostrar
resultats i mostrar errors. \texttt{SimulationViewListener} és el contracte invers:
la vista l'empra per avisar que l'usuari ha demanat una simulació.

\texttt{SimulationFrame} és la finestra Swing concreta. Permet indicar $N$, una
llavor opcional, la variant de regles, el nombre de jugadors, el nombre de fils i si
es vol comparar variants. Els resultats es mostren en cinc formes: estadístics
principals, text amb el format exacte de l'enunciat, histograma de torns, taula de
freqüències de visita i taula de probabilitats de victòria.

L'histograma no agrupa els torns en intervals continus: cada barra correspon a un
valor enter de torns observats. Això és coherent amb la naturalesa discreta de la
variable $X$. Els eixos es dibuixen amb marques regulars de valors ``bons''
(múltiples de 1, 2, 5 i potències de 10), de manera que l'escala sigui estable i
llegible encara que el màxim observat canviï entre simulacions. Aquesta lògica de
dibuix queda encapsulada a \texttt{HistogramPanel}, i per això \texttt{SimulationFrame}
no necessita conèixer els detalls de pintat de l'eix ni de les barres.

\subsection{CLI Reproduïble}

La classe \texttt{CLI} és l'eina emprada per generar les taules d'aquesta memòria.
Permet passar valors de $N$, llavor base, variants, repeticions, nombre de
jugadors, nombre de fils, format de sortida i nombre de caselles més visitades. La
Llista~\ref{lst:ordres} mostra un exemple representatiu amb sortida CSV abreujada.

\begin{lstlisting}[caption={Exemple d'execució reproduïble de la CLI.},label={lst:ordres}]
CLI --games 100000 --seed 20260613 --players 4 \
  --variants standard --repetitions 1 --format csv

variant,players,n,rep,...,mean,win1,...,win4
STANDARD,4,...,45.831060,0.262760,...,0.237110
\end{lstlisting}

La CLI imprimeix la llavor efectiva de cada fila. Això és important perquè les
llavors de les repeticions no són simplement consecutives: es deriven de la llavor
base, de la variant, del valor de $N$ i del número de repetició. Així, una fila es
pot repetir exactament sense dependre de les altres.

%% ================================================================
\section{Metodologia Experimental}
\label{sec:experiments}

Les conclusions estadístiques provenen de les mostres de torns i de les llavors
impreses per la CLI. Per evitar barrejar efectes de paral·lelització amb les
estimacions principals, les taules estadístiques s'han generat amb un sol fil. La
concurrència s'analitza en una prova separada.

S'han dissenyat cinc grups de proves:
\begin{enumerate}
    \item \textbf{Convergència:} 5 repeticions independents per a $N=1\,000$,
    $10\,000$, $100\,000$, $250\,000$, $500\,000$ i $1\,000\,000$ amb regles
    oficials.
    \item \textbf{Variants:} 3 repeticions amb $N=250\,000$ per a cada variant.
    \item \textbf{Jugadors:} 3 repeticions amb $N=250\,000$ per a 1, 2, 3 i 4
    jugadors amb regles oficials.
    \item \textbf{Freqüències de visita:} una execució de $N=250\,000$ amb regles
    oficials i les 12 caselles més visitades.
    \item \textbf{Concurrència:} $N=1\,000\,000$ amb 1, 2, 4 i 8 fils.
\end{enumerate}

A partir de la prova de convergència s'ha triat $N^*=250\,000$ per a variants,
jugadors i visites. Aquest valor manté mediana i percentils estables i redueix
l'error estàndard de la mitjana a prop de 0.013 torns; els increments fins a
$500\,000$ i $1\,000\,000$ milloren la precisió, però no canvien les conclusions
substantives. La prova de concurrència conserva $N=1\,000\,000$ perquè necessita un
lot gran per observar diferències entre nombres de fils.

Quan una taula agrega diverses repeticions, la columna de mitjana indica la mitjana
de les mitjanes obtingudes en aquelles repeticions. Quan es mostra una execució
individual, la mitjana correspon només als torns d'aquella execució. Els percentils
provenen del criteri de rang més proper de l'Equació~\ref{eq:percentil}.

%% ================================================================
\section{Resultats i Discussió}

\subsection{Convergència de la Mitjana i dels Percentils}

La Taula~\ref{tab:convergencia} resumeix les 30 execucions de convergència. Per
cada valor de $N$ s'ha calculat la mitjana de les cinc mitjanes obtingudes, la
desviació típica mostral d'aquestes mitjanes, el rang de mitjanes observat, la
mediana típica i el percentil $0.9$ típic.

\begin{table*}[t]
\centering
\caption{Convergència amb regles oficials. Cada fila agrega 5 repeticions independents.}
\label{tab:convergencia}
\footnotesize
\begin{tabular}{rrrrrrr}
\toprule
$N$ & Mitjana agregada & Desv. típica & Rang de mitjanes & Mediana & $P_{0.1}$ & $P_{0.9}$ \\
\midrule
1\,000 & 20.757800 & 0.4008 & 20.070000--21.051000 & 16--17 & 8 & 37--39 \\
10\,000 & 20.714960 & 0.1460 & 20.453900--20.787000 & 17 & 8 & 38--39 \\
100\,000 & 20.769808 & 0.0352 & 20.729070--20.807100 & 17 & 8 & 38--39 \\
250\,000 & 20.759382 & 0.0282 & 20.731120--20.805856 & 17 & 8 & 38--39 \\
500\,000 & 20.753412 & 0.0262 & 20.719298--20.783472 & 17 & 8 & 38--39 \\
1\,000\,000 & 20.761862 & 0.0119 & 20.749209--20.776132 & 17 & 8 & 39 \\
\bottomrule
\end{tabular}
\end{table*}

La reducció de variabilitat és clara. Amb $N=1\,000$, el rang de mitjanes és de
0.981 torns. Amb $N=100\,000$, baixa a 0.078 torns, i amb $N=1\,000\,000$ queda en
0.027 torns. Això concorda amb el comportament esperat a partir de
l'equació~\ref{eq:error-mitjana}: l'error típic de la mitjana decreix de manera
aproximada com $O(1/\sqrt{N})$.

\begin{figure}[H]
\centering
\begin{tikzpicture}
\begin{semilogxaxis}[
    width=\columnwidth,
    height=0.58\columnwidth,
    xlabel={$N$},
    ylabel={Desv. típica de la mitjana},
    xmin=800, xmax=1200000,
    ymin=0, ymax=0.43,
    xtick={1000,10000,100000,250000,500000,1000000},
    xticklabels={$10^3$,$10^4$,$10^5$,$2.5\cdot10^5$,$5\cdot10^5$,$10^6$},
    ymajorgrids=true,
    grid style={gray!25},
    tick label style={font=\scriptsize},
    label style={font=\scriptsize}
]
\addplot+[mark=*] coordinates {
    (1000,0.4008)
    (10000,0.1460)
    (100000,0.0352)
    (250000,0.0282)
    (500000,0.0262)
    (1000000,0.0119)
};
\end{semilogxaxis}
\end{tikzpicture}
\caption{Reducció de la variabilitat de la mitjana estimada quan creix $N$.}
\label{fig:convergencia-sd}
\end{figure}

Els percentils són encara més estables. El percentil $0.1$ és 8 en totes les
execucions de la taula, la mediana queda pràcticament fixada en 17 i el percentil
$0.9$ oscil·la només entre 38 i 39 per a $N\geq10\,000$. Això suggereix que la
forma central de la distribució es pot descriure amb menys mostres que les que fan
falta per estabilitzar els extrems de la cua.

\subsection{Incertesa de l'Estimació}

Les cinc repeticions per cada valor de $N$ també permeten estimar la incertesa de
la mitjana estimada. Si $s$ és la desviació típica mostral de les cinc mitjanes,
l'error estàndard de la mitjana d'aquestes cinc execucions és $s/\sqrt{5}$. Per a
$N=1\,000$, és aproximadament 0.179 torns; per a $N=100\,000$, baixa a 0.016 torns;
per a $N=250\,000$, és 0.013 torns; i per a $N=1\,000\,000$, és 0.005 torns.

Aquesta anàlisi no pretén donar un interval de confiança formal molt precís, perquè
cinc repeticions són poques per estimar una variància amb alta fiabilitat. Tot i
això, és suficient per justificar la decisió experimental: s'utilitza
$N^*=250\,000$ per a variants, multijugador i visites. Aquest valor ja mostra una
mitjana estable i percentils centrals pràcticament constants. Les mostres de
$500\,000$ i $1\,000\,000$ milloren la precisió, però no canvien la interpretació
de les comparacions principals.

Una altra observació important és que el mínim observat cau de 3 a 1 quan $N$
augmenta. Això no indica un error: una partida pot acabar en un sol torn si, dins
el mateix torn, diverses tirades d'oca encadenen salts fins a la casella final. A
mostres petites és poc probable veure aquest cas extrem, però amb $100\,000$
partides ja apareix en diverses repeticions.

\subsection{Cua de la Distribució}

La part menys estable és el màxim. En les cinc execucions amb $N=1\,000$, els màxims
observats foren 100, 128, 93, 89 i 95. Amb $N=100\,000$, apareixen màxims de fins a
201, i amb $N=1\,000\,000$ apareix una execució amb 233 torns. Aquest creixement no
contradiu la convergència de la mitjana: el màxim mesura un extrem de la cua, no el
comportament central.

La cua llarga és coherent amb les regles. Caselles com la mort poden retornar el
jugador al principi, i les penalitzacions poden afegir torns sense progrés. Això
crea partides molt llargues amb probabilitat baixa. Per aquesta raó, el màxim és
útil com a observació de robustesa, però no s'ha d'interpretar com una estimació
estable del ``cas pitjor'' real.

\subsection{Comparació de Variants}

La Taula~\ref{tab:variants} compara variants de regles amb $N=250\,000$ i tres
repeticions. Per condensar la informació, es mostra la mitjana de les tres
mitjanes, la mediana típica, el percentil $0.9$ i la diferència respecte de les
regles oficials.

\begin{table*}[t]
\centering
\caption{Efecte de les variants de regles amb $N=250\,000$ i 3 repeticions.}
\label{tab:variants}
\footnotesize
\begin{tabular}{lrrrrr}
\toprule
Variant & Mitjana agregada & Mediana & $P_{0.9}$ & Diferència & Diferència \% \\
\midrule
Regles oficials & 20.753311 & 17 & 38--39 & 0.000000 & 0.0 \\
Sense penalitzacions & 18.695409 & 15 & 34 & -2.057901 & -9.9 \\
Oques sense tirada extra & 24.998955 & 20 & 46 & +4.245644 & +20.5 \\
Sense caselles especials & 22.771131 & 21 & 30 & +2.017820 & +9.7 \\
\bottomrule
\end{tabular}
\end{table*}

La variant sense penalitzacions és la més ràpida: redueix la mitjana en uns dos
torns, gairebé un 10\%. Això confirma que posada, pou i presó tenen un impacte
important en la durada, sobretot perquè afegeixen torns sense cap moviment.

La variant sense tirada extra d'oca és la més lenta. Les oques continuen
teletransportant cap endavant, però l'avantatge de no comptar la nova tirada
desapareix. La mitjana puja més de quatre torns. Això mostra que la regla ``de oca
a oca y tiro porque me toca'' no és un detall menor: és una de les fonts principals
de reducció del nombre de torns comptats.

La variant sense caselles especials té una mitjana superior a l'oficial, però un
percentil $0.9$ menor. Això sembla paradoxal només al principi. Sense caselles
especials no hi ha dreceres d'oca ni salts endavant, de manera que la partida típica
triga més. Però tampoc hi ha mort, pou ni presó, de manera que la cua dreta es fa
menys extrema. El resultat és una distribució més concentrada però desplaçada cap a
valors centrals una mica més alts.

\subsection{Partides amb Diversos Jugadors}

La Taula~\ref{tab:jugadors} mostra l'extensió opcional a partides de 1 a 4
jugadors. En aquesta prova tots els jugadors segueixen les mateixes regles i juguen
en ordre cíclic. Quan una fitxa acaba el torn en una casella ocupada, captura
l'ocupant anterior i l'envia a la casella 0. Per això la mitjana representa el
nombre total de torns de la partida col·lectiva, no els torns consumits per cada
jugador individualment. Cada fila agrega tres repeticions de $250\,000$ partides,
és a dir, $750\,000$ partides per cada nombre de jugadors.

\begin{table*}[t]
\centering
\caption{Durada i probabilitats de victòria amb 1--4 jugadors, $N=250\,000$ i 3 repeticions.}
\label{tab:jugadors}
\footnotesize
\begin{tabular}{rrrrrrrr}
\toprule
Jugadors & Mitjana agregada & Mediana & $P_{0.9}$ & $\Pr(W=1)$ & $\Pr(W=2)$ & $\Pr(W=3)$ & $\Pr(W=4)$ \\
\midrule
1 & 20.753311 & 17 & 38--39 & 1.000000 & -- & -- & -- \\
2 & 29.249520 & 26 & 50 & 0.510141 & 0.489859 & -- & -- \\
3 & 37.556737 & 34 & 62 & 0.345716 & 0.333723 & 0.320561 & -- \\
4 & 45.791199 & 42 & 74 & 0.263187 & 0.254387 & 0.244127 & 0.238300 \\
\bottomrule
\end{tabular}
\end{table*}

La Taula~\ref{tab:jugadors-biaix} aprofundeix en la mateixa prova. El rang de
mitjanes és el mínim i el màxim de les tres repeticions. La columna
mitjana/$J$ no és una durada individual completa, sinó una normalització del
cost total de la partida per nombre de jugadors; serveix per veure com el fet
d'aturar-se al primer guanyador redueix el cost efectiu per jugador.

\begin{table*}[t]
\centering
\caption{Estabilitat i biaix de l'ordre de torn en l'experiment multijugador.}
\label{tab:jugadors-biaix}
\footnotesize
\begin{tabular}{rrrrrrrr}
\toprule
Jugadors & Rang de mitjanes & Mitjana/$J$ & $1/J$ & $\Pr(W=1)-1/J$ & $\Pr(W=J)-1/J$ & Biaix $1-J$ & Ràtio $W_1/W_J$ \\
\midrule
1 & 20.744856--20.757584 & 20.753311 & 1.000000 & -- & -- & -- & -- \\
2 & 29.217540--29.274840 & 14.624760 & 0.500000 & +0.010141 & -0.010141 & 0.020283 & 1.041 \\
3 & 37.526280--37.590728 & 12.518912 & 0.333333 & +0.012383 & -0.012772 & 0.025155 & 1.078 \\
4 & 45.773976--45.815492 & 11.447800 & 0.250000 & +0.013187 & -0.011700 & 0.024887 & 1.104 \\
\bottomrule
\end{tabular}
\end{table*}

\begin{figure}[H]
\centering
\begin{tikzpicture}
\begin{axis}[
    width=\columnwidth,
    height=0.62\columnwidth,
    xlabel={Jugador},
    ylabel={Probabilitat de victòria},
    xmin=0.8, xmax=4.2,
    ymin=0.20, ymax=0.54,
    xtick={1,2,3,4},
    ymajorgrids=true,
    grid style={gray!25},
    legend style={font=\scriptsize, at={(0.98,0.98)}, anchor=north east},
    tick label style={font=\scriptsize},
    label style={font=\scriptsize}
]
\addplot+[mark=*] coordinates {(1,0.510141) (2,0.489859)};
\addlegendentry{2 jugadors}
\addplot+[mark=square*] coordinates {(1,0.345716) (2,0.333723) (3,0.320561)};
\addlegendentry{3 jugadors}
\addplot+[mark=triangle*] coordinates {(1,0.263187) (2,0.254387) (3,0.244127) (4,0.238300)};
\addlegendentry{4 jugadors}
\end{axis}
\end{tikzpicture}
\caption{Probabilitat de victòria segons l'ordre de torn amb captures.}
\label{fig:victories}
\end{figure}

La durada total creix amb el nombre de jugadors, però no de manera proporcional. Amb
quatre jugadors la mitjana és 45.79 torns totals, molt per sota de quatre partides
individuals independents, perquè només cal que un jugador acabi. La mediana puja de
17 a 42 torns totals i el percentil $0.9$ arriba a 74, cosa que indica que la
partida col·lectiva continua tenint cua llarga. Ara bé, la columna
La mitjana dividida per $J$ baixa de 20.75 torns amb un jugador a 11.45 amb quatre jugadors:
afegir jugadors incrementa la durada total i les captures introdueixen retrocessos,
però també augmenta la probabilitat que algun jugador tengui una trajectòria ràpida.

Les repeticions són estables. El rang de mitjanes és de 0.013 torns amb un jugador,
0.057 amb dos, 0.064 amb tres i 0.042 amb quatre. Aquestes diferències són petites
comparades amb els canvis entre nombres de jugadors, que són de 8 a 9 torns totals
entre files consecutives. Per tant, les conclusions sobre durada no depenen d'una
llavor concreta.

Les probabilitats de victòria no són exactament uniformes. Si l'ordre de torn no
introduís cap biaix, cada jugador hauria de guanyar amb probabilitat $1/J$. En
canvi, el primer jugador guanya un 51.0\% de les partides de dos jugadors i un
26.3\% de les de quatre, mentre que el darrer baixa fins al 23.8\% en el cas de
quatre. El biaix entre el primer i el darrer jugador es manté al voltant de 2.5
punts percentuals; en termes relatius, el primer jugador guanya aproximadament un
10.7\% més de partides que el quart.

Aquest efecte és massa gran per atribuir-lo al soroll Monte Carlo. Amb $750\,000$
partides agregades, l'error estàndard d'una proporció al voltant de $0.25$ és de
l'ordre de $5\cdot10^{-4}$, molt menor que una desviació d'1.0--1.3 punts
percentuals respecte de $1/J$. El mecanisme és estructural: els torns són cíclics i
la partida acaba immediatament quan algú arriba a la casella 63. Si el primer
jugador acaba en una ronda, els jugadors posteriors ja no reben l'oportunitat
d'igualar-lo en aquella mateixa ronda. Si es volgués una experiència simètrica,
bastaria aleatoritzar quin jugador comença cada partida.

\subsection{Freqüències de Visita}

La Taula~\ref{tab:visites} mostra les 12 caselles més visitades en una execució de
$N=250\,000$ amb regles oficials. El percentatge està calculat sobre el total de
visites registrades, no sobre el nombre de partides. Una partida pot visitar moltes
caselles, i una mateixa casella pot aparèixer més d'una vegada en la mateixa
partida.

\begin{table}[H]
\centering
\caption{Caselles més visitades amb regles oficials, $N=250\,000$.}
\label{tab:visites}
\footnotesize
\begin{tabular}{rrrr}
\toprule
Rang & Casella & Visites & Percentatge \\
\midrule
1 & 0 & 364\,644 & 4.4179 \\
2 & 52 & 315\,752 & 3.8256 \\
3 & 31 & 291\,456 & 3.5312 \\
4 & 63 & 250\,000 & 3.0289 \\
5 & 41 & 213\,077 & 2.5816 \\
6 & 59 & 212\,297 & 2.5721 \\
7 & 36 & 211\,987 & 2.5684 \\
8 & 62 & 201\,854 & 2.4456 \\
9 & 6 & 187\,316 & 2.2695 \\
10 & 12 & 187\,316 & 2.2695 \\
11 & 61 & 184\,460 & 2.2349 \\
12 & 30 & 181\,992 & 2.2050 \\
\bottomrule
\end{tabular}
\end{table}

La casella 0 és la més visitada perquè totes les partides hi comencen i perquè la
casella de mort 58 retorna el jugador al principi. Les caselles 52 i 31 també són
molt freqüents perquè la visita es compta tant quan s'hi cau com durant els torns
anul·lats que mantenen el jugador immòbil. La casella 63 apareix exactament
250\,000 vegades: una per partida acabada.

Les caselles 41, 59 i 36 formen part de la cadena d'oques. La seva presència
entre les més visitades és esperable perquè una caiguda en una oca pot generar més
tirades dins el mateix torn i travessar diverses oques en una seqüència curta. La
simetria entre 6 i 12 prové del pont: quan una d'aquestes caselles apareix, l'altra
és la destinació immediata. La casella 30 també apareix entre les primeres perquè és
la destinació del laberint.

\subsection{Concurrència}

La Taula~\ref{tab:threads} mostra una prova de rendiment amb $N=1\,000\,000$ i la
mateixa llavor base. Com que la simulació paral·lela reparteix la llavor entre
treballadors, les mostres no són idèntiques entre nombres de fils diferents, però
sí són reproduïbles per a cada configuració concreta.

\begin{table}[H]
\centering
\caption{Temps de simulació amb $N=1\,000\,000$ i regles oficials.}
\label{tab:threads}
\footnotesize
\begin{tabular}{rrrrr}
\toprule
Fils & Temps (ms) & Acceleració & Mitjana & Màxim \\
\midrule
1 & 573 & 1.00 & 20.759707 & 233 \\
2 & 364 & 1.57 & 20.767541 & 212 \\
4 & 176 & 3.26 & 20.765076 & 183 \\
8 & 121 & 4.74 & 20.755403 & 183 \\
\bottomrule
\end{tabular}
\end{table}

El guany amb 8 fils és d'aproximadament $4.7\times$, lluny d'un factor 8 però
clarament positiu. Hi ha diverses raons: la creació i coordinació de fils no és
gratuïta, cal fusionar vectors de freqüències i la part d'ordenació dels $N$ torns
continua sent seqüencial en aquesta implementació. A més, per a $N$ relativament
petits la sobrecàrrega pot dominar. Per això la GUI permet configurar el nombre de
fils en lloc d'imposar sempre el màxim disponible.

%% ================================================================
\section{Conclusions}

La pràctica s'ha resolt amb una simulació Monte Carlo modular i extensible. El
model de joc queda concentrat a \texttt{Game}, la simulació massiva a
\texttt{Simulation} i els estadístics a \texttt{SimulationStats}. La sortida
obligatòria de \texttt{Main} és simple i estricta, mentre que \texttt{CLI} ofereix
les opcions necessàries per fer experiments reproduïbles sense interacció manual.
La GUI compleix el requisit d'interfície gràfica i afegeix histograma discret amb
eixos llegibles, comparació de variants, selecció de jugadors, probabilitats de
victòria i freqüències de visita.

Els resultats empírics indiquen que la durada típica d'una partida amb regles
oficials és d'aproximadament 20.75 torns, amb mediana 17 i percentil $0.9$ al
voltant de 38--39. La mitjana convergeix de manera visible quan $N$ creix: la
desviació entre repeticions baixa de 0.4008 torns amb $N=1\,000$ a 0.0282 torns amb
$N=250\,000$ i a 0.0119 amb $N=1\,000\,000$. En canvi, el màxim observat és
inestable i creix amb la mida de la mostra, fet coherent amb una distribució de cua
llarga.

L'anàlisi de variants mostra que les penalitzacions afegeixen prop d'un 10\% de
durada, mentre que eliminar la tirada extra de les oques incrementa la mitjana més
d'un 20\%. Les partides multijugador amb captures confirmen un avantatge petit però
persistent del primer jugador: amb quatre jugadors guanya aproximadament el 26.3\%
de les partides, mentre que el quart queda al 23.8\%. Les captures també allarguen
les partides col·lectives perquè poden retornar fitxes avançades a la sortida. Les
freqüències de visita confirmen el paper especial de la mort, el pou i la presó, i
també mostren la presència destacada de la cadena d'oques. Finalment, la simulació
concurrent millora el temps d'execució per lots grans, però amb guanys sublineals
per la sobrecàrrega i per la part seqüencial del procés estadístic.

Com a treball futur, seria natural exportar totes les freqüències de visita a CSV,
comparar altres convencions de percentils i calcular la distribució exacta
mitjançant una cadena de Markov per al cas d'un jugador. També es podria afegir una
opció per aleatoritzar el jugador inicial i comprovar que les probabilitats de
victòria es tornen simètriques.

\end{document}
